\section{Architecture: Proof Outsourcing}
\label{sec:design}

The previous section establishes that supercompilation \emph{is} cyclic proof
search.  This section makes it operational: the proof engine (supercompiler)
delegates creative steps to an untrusted oracle (LLM), and the kernel validates
every proposal deterministically.

\subsection{The proof engine}

The proof engine is the supercompiler's main loop: it receives a judgement
$\jTy{\Gamma}{t}{A}$ and constructs a cyclic proof graph.  At each vertex:
\begin{itemize}[leftmargin=*]
\item \textbf{Invertible phase:} Apply driving rules ($\beta$/$\iota$/$\delta$) deterministically.
\item \textbf{Synchronous phase:} When facing a neutral case-scrutinee, split
  into one successor per compatible constructor.
\item \textbf{Cyclic closure:} When a configuration matches a previously-seen
  one, close the graph with a back-edge.
\end{itemize}

When neither driving nor splitting applies, the engine needs a \emph{cut} --- a
generalisation that makes folding possible.  This is where the oracle enters.

\subsection{The untrusted oracle}

The oracle is a function $\texttt{config} \to \texttt{list (config * nat)}
\to \texttt{option gen\_result}$ implemented by an LLM.  Given the current
stuck configuration and candidates from the memo table, it proposes a
generalised configuration and instantiations.  The proposal is validated:
\begin{itemize}[leftmargin=*]
\item \textbf{Instance check:} Both the current and companion configs must be
  instances of the proposal (substitution check).
\item \textbf{Trace check:} The resulting proof graph must pass
  \texttt{trace\_condition\_ok} (no ungrounded cycles).
\end{itemize}

The validation is deterministic; the oracle is untrusted for soundness.
If the LLM proposes a valid generalisation, the proof proceeds.  If it
proposes nonsense, the kernel rejects it and the engine continues driving.

\subsection{The $\omega$-rule: lemma outsourcing}

When the proof engine is stuck and the oracle cannot propose a useful
generalisation, a \emph{lemma} may unblock it.  The $\omega$-rule outsources
lemma synthesis:

\begin{enumerate}[leftmargin=*]
\item The LLM proposes a lemma statement (a pair of terms $\mathit{lhs},
  \mathit{rhs}$).
\item The kernel runs the proof engine on both $\mathit{lhs}$ and $\mathit{rhs}$
  independently.  If both produce identical residuals, the lemma is validated.
\item The validated lemma is added to the proof engine as a rewrite rule.
\item The main proof retries with the lemma, which may enable a fold that was
  previously blocked.
\end{enumerate}

This is the same pattern at two levels: the LLM proposes cuts for the main
proof (generalisation) and for auxiliary lemmas ($\omega$-rule).  The kernel
validates both through the same proof engine.

\subsection{The pipeline}

\begin{center}
\begin{tikzpicture}[
  box/.style={rectangle,draw,rounded corners,fill=blue!5,
              text width=4cm,align=center,font=\small,minimum height=0.8cm},
  trybox/.style={box,fill=green!8},
  arrow/.style={->,>=stealth,thick},
  node distance=1.0cm
]
\node[box,fill=red!8] (stuck) {Config stuck\\(AU returns None)};
\node[trybox]        (llm)   [below=of stuck] {LLM generalisation\\(cut introduction)};
\node[trybox]        (lem)   [below=of llm]   {LLM lemma proposal\\($\omega$-rule)};
\node[box,draw=none,fill=none] (fail) [below=of lem] {drive forward\\(no fold)};

\draw[arrow] (stuck) -- node[right,font=\scriptsize]{try} (llm);
\draw[arrow] (llm)   -- node[right,font=\scriptsize]{None} (lem);
\draw[arrow] (lem)   -- node[right,font=\scriptsize]{None} (fail);
\end{tikzpicture}
\end{center}

When the pipeline succeeds, the LLM is transparent: the resulting proof is
a valid cyclic derivation checked by the kernel.  The LLM is a heuristic for
proof search, not a trusted component.
